% !TeX root = main.tex
\section{Dynamics Theory}

\subsection{Purpose and Authority}

This section defines the approved physical and numerical theory shared by the
active particle dynamics.  Model source documents describe implementation;
\texttt{ProblemEvolution.tex} records historical investigations.  Retired
birth schedules, effective radii, manual-wall variants, and visualization
choices are not dynamics theory and do not belong here.

The theory is momentum governed.  Energy is not a governing simulation state
and is never used to repair motion after a frame.  Optional energy calculations
are diagnostics only.

\subsection{Core Invariants}

Every implementation shall preserve the following invariants:

\begin{enumerate}
    \item Physical particle radii always define contact geometry.
    \item Physical overlap always creates a physical contact.
    \item A contact ends only after the physical objects separate.
    \item Contact history may affect stored local state, but it may not alter
          geometry or make overlapping particles transparent.
    \item Every source reads one immutable frame-start state.
    \item A source writes only its own velocity, position, contact state, and
          lifecycle state.
    \item All contacts affecting a source are known before that source's final
          velocity is accepted.
    \item Boundary markers provide locality; analytic geometry defines walls.
    \item Numerical containment may block motion, but blocked momentum must be
          accounted for explicitly.
\end{enumerate}

Initial particle-particle or particle-wall overlap is invalid input.  It is
reported before the first dynamics frame rather than reinterpreted through
reduced radii or zero-force starting-contact state.

\subsection{Particle State}

A mobile particle has position \(\mathbf x_i\), velocity \(\mathbf v_i\),
mass \(m_i\), and physical radius \(r_i\).  Its translational momentum is

\begin{equation}
\mathbf p_i=m_i\mathbf v_i.
\end{equation}

Mass and radius are positive.  Position and velocity are double buffered in
the Vulkan representation.  One selector identifies the immutable current
buffers; integration writes the alternate buffers.  Python uses an equivalent
frame snapshot so contact results do not depend on source iteration order.

Particle index zero is reserved as the null particle.  Boundary-marker
particles are geometric locality records, not mobile sources.  Dead particles
are excluded from source processing, contact targets, cell occupancy, and
rendering.

\subsection{Source-Owned Execution}

One dynamics frame follows this conceptual order:

\begin{enumerate}
    \item snapshot current positions and velocities;
    \item validate initial geometry when the configuration is first used;
    \item for each active source, discover every current particle, piston, and
          stationary-wall contact;
    \item accumulate one source-local force sum;
    \item calculate the source's candidate velocity;
    \item apply the source's complete maximum-depth constraint set;
    \item integrate the source position with the accepted velocity;
    \item apply cell-domain and death-bound rules;
    \item swap the current and alternate state buffers.
\end{enumerate}

A source never writes a target particle.  For a mobile pair, each source
invocation independently evaluates the same frame-start geometry.  A symmetric
pair law must therefore produce equal and opposite impulses without relying on
invocation order.

\subsection{Physical Particle Contact}

For source \(i\) and target \(j\), define

\begin{align}
\mathbf d_{ij} &= \mathbf x_j-\mathbf x_i,\\
d_{ij} &= \lVert\mathbf d_{ij}\rVert,\\
\mathbf n_{ij} &= \frac{\mathbf d_{ij}}{d_{ij}}.
\end{align}

The normal points from source to target.  Coincident centers use a fixed
deterministic fallback normal.  Physical penetration is

\begin{equation}
\delta_{ij}=\max(0,r_i+r_j-d_{ij}).
\end{equation}

A contact exists exactly when \(d_{ij}<r_i+r_j\).  Geometry is calculated
once and passed through contact initialization, force accumulation, and
constraint construction.  Velocity does not change instantaneous overlap.

The optional circular overlap-area measure is denoted \(A_{ij}\).  Both
\(\delta_{ij}\) and \(A_{ij}\) are symmetric geometric quantities.

\subsection{Contact Force}

The current contact force magnitude is

\begin{equation}
F_{ij}=q_{ij}M_{ij},
\end{equation}

where \(q_{ij}\geq0\) is symmetric pair stiffness and \(M_{ij}\) is the
configured contact measure:

\begin{equation}
M_{ij}=\begin{cases}
\delta_{ij}, & \text{depth mode},\\
A_{ij}, & \text{area mode}.
\end{cases}
\end{equation}

The force on the source is repulsive:

\begin{equation}
\mathbf F_{i\leftarrow j}=-F_{ij}\mathbf n_{ij}.
\end{equation}

The frame impulse generated by this force is

\begin{equation}
\mathbf J_{i\leftarrow j}=\mathbf F_{i\leftarrow j}\Delta t.
\end{equation}

This ordinary force impulse changes translational momentum directly.  It must
not also be recorded as convertible internal momentum.

The approved material direction is to replace constant stiffness with a
reversible depth-dependent stiffness while keeping the same momentum-governed
state.  Energy is not added as a dynamic variable.  Instead, the contact force
remains a deterministic function of current physical penetration:
\[
    x=\mathrm{clamp}
      \left(\frac{\delta}{\delta_{\mathrm{hard}}},0,1\right),
\]
\[
    q_{\mathrm{eff}} =
      q_{\mathrm{base}}\left(1+g x^p\right),
\]
\[
    F=q_{\mathrm{eff}}M.
\]
In depth mode \(M=\delta\).  The same curve is used during compression and
release, so the impulse produced by ordinary force integration slows, stops,
and reverses particles without a separate energy ledger.  For
particle-particle contacts, the stiffness law must remain symmetric so the
two source-owned evaluations produce equal and opposite force directions from
the same frame-start geometry.

\subsection{Parametric Wall Theory}

Every stationary wall is a finite parametric curve segment.  In two
dimensions, a cubic segment is

\begin{align}
x(t)&=a_x+b_xt+c_xt^2+d_xt^3,\\
y(t)&=a_y+b_yt+c_yt^2+d_yt^3,
\qquad 0\leq t\leq1.
\end{align}

Straight walls are the same representation with zero quadratic and cubic
coefficients.  The closest bounded point on the segment supplies the tangent.
The configured orientation supplies the outward normal.  Contact distance and
penetration are calculated from that analytic point and normal, never from the
orientation of a marker particle.

Boundary particles with positive \texttt{ptype} are broad-phase locality
markers.  The generic simple model no longer stores a wall evaluator in
\texttt{Data.z}; every stationary wall is described by
\texttt{curve\_wall\_segments}.  If any marker is local to a source, all
applicable configured curve segments are evaluated.  At most one contact is
retained for each positive physical \texttt{wall\_flag}: the deepest valid
contact.  Marker particle IDs are not physical wall identities.
Generator-only curve roles may identify packing curves and boundary curves,
but those roles are not runtime evaluator selectors.  Runtime wall response is
defined by the configured curve geometry and physical wall flag.

A stationary wall is an external momentum sink or source.  Particle-only
momentum is therefore not a closed-system invariant when stationary walls are
active.  The equal-and-opposite reaction belongs to the wall or environment,
even when that external body is not integrated.

\subsection{Moving Piston Theory}

The piston is an analytic moving plane, not a row of moving boundary markers.
Before \texttt{piston\_start\_frame}, it is stationary at the chamber start.
It then advances with the configured piston velocity and clamps at the chamber
end, where its velocity becomes zero.

Each active source evaluates the piston once, independently of boundary-marker
locality.  Piston contact uses velocity relative to the moving plane for
timestep resolution and maximum-depth limits.  The particle impulse has an
equal-and-opposite piston reaction:

\begin{equation}
\Delta\mathbf p_{\mathrm{piston}}
=-\Delta\mathbf p_{\mathrm{particle}}.
\end{equation}

Particle momentum plus accumulated piston reaction is the appropriate
momentum diagnostic for a piston-driven chamber.

\subsection{Velocity and Position Integration}

After the complete source force is known, semi-implicit Euler calculates

\begin{align}
\mathbf v_i^* &= \mathbf v_i^n+
    \frac{\mathbf F_i}{m_i}\Delta t,\\
\mathbf x_i^{n+1} &= \mathbf x_i^n+
    \mathbf v_i^{n+1}\Delta t.
\end{align}

Here \(\mathbf v_i^*\) is the force-driven candidate.  Maximum-depth
containment may replace it with the accepted velocity
\(\mathbf v_i^{n+1}\).  Position always uses the accepted new velocity.

\subsection{Maximum Penetration}

The model limits observed penetration to a fraction of the source radius:

\begin{equation}
\delta_{\mathrm{target}}=f_{\mathrm{target}}r_i,
\qquad
\delta_{\mathrm{hard}}=f_{\mathrm{hard}}r_i,
\qquad
0<f_{\mathrm{target}}<f_{\mathrm{hard}}<1.
\end{equation}

The target depth is the preferred containment depth.  Below this depth the
ordinary contact force may continue to compress the contact.  At or beyond
this depth, additional inward normal motion is blocked and the blocked
momentum is stored as convertible internal momentum.  The hard depth is the
fatal observed compression limit.  Penetration beyond the hard depth is an
error.

For contact normal \(\mathbf n_k\) and target normal velocity \(u_k\), a
contact at or beyond target depth contributes the unilateral velocity
constraint

\begin{equation}
\mathbf v_i^{n+1}\cdot\mathbf n_k\leq u_k.
\end{equation}

All active maximum-depth contacts constrain the same candidate velocity.
The two-dimensional solver merges nearly parallel constraints, retains the
stricter limit, and solves no more than two independent normals.  Non-finite,
singular, or unsatisfied projections are errors; the candidate velocity must
not be applied after solver failure.

The momentum blocked by containment is

\begin{equation}
\Delta\mathbf p_{\mathrm{blocked}}
=m_i\left(\mathbf v_i^*-\mathbf v_i^{n+1}\right).
\end{equation}

This quantity is not ordinary force impulse.  It is the source of convertible
internal momentum defined below.

\subsection{Timestep Resolution and Tunneling}

Containment is valid only when one timestep cannot cross the unused
penetration reserve.  With inward relative normal speed
\(v_{\mathrm{inward}}\), require

\begin{equation}
v_{\mathrm{inward}}\Delta t
\leq\delta_{\mathrm{hard}}-\delta.
\end{equation}

Violation is a timestep-resolution error and requires a smaller
\(\Delta t\); containment is not a substitute for unresolved temporal
sampling.

Generator-side feasibility checks should report the same scale before the
simulation runs.  For a closing speed \(v_c\), the per-frame closing distance
is \(v_c\Delta t\).  The useful configuration ratios are
\[
    \frac{v_c\Delta t}{\delta_{\mathrm{target}}}
    \qquad\mathrm{and}\qquad
    \frac{v_c\Delta t}{\delta_{\mathrm{hard}}}.
\]
Values near or above one mean that a single frame can consume the available
compression depth.  Increasing particle radius, reducing closing speed,
reducing \(\Delta t\), or increasing the allowed penetration fractions changes
this numerical feasibility.  The feasibility check does not replace the
runtime max-depth rule; it explains whether a chosen configuration is likely
to enter overcompression.

Current-state tunneling is also an error.  For particle contact, a target's
leading perimeter may not pass the source center.  For wall contact, the wall
ghost's leading perimeter may not pass the particle center.  A later swept
contact test may detect complete between-frame crossings, but it must not
replace these observed-state checks.

\subsection{Momentum Accounting}

\subsubsection{Translational Momentum}

Translational momentum is represented only by particle velocity:

\begin{equation}
\mathbf p_{\mathrm{trans},i}=m_i\mathbf v_i.
\end{equation}

An ordinary contact impulse belongs here.  Storing the same impulse in another
ledger would count it twice.

\subsubsection{Convertible Internal Momentum}

Convertible internal momentum is momentum temporarily removed from
translational motion by a contact constraint and available for later return to
translational motion.  For each active contact \(c\), store a vector

\begin{equation}
\mathbf p_{\mathrm{conv},i,c}.
\end{equation}

Only blocked constraint momentum may increase this value.  It is owned per
contact because simultaneous contacts can have different normals, targets,
and release conditions.  The particle's momentum accounting is

\begin{equation}
\mathbf p_{\mathrm{accounted},i}
=m_i\mathbf v_i+
\sum_c\mathbf p_{\mathrm{conv},i,c}.
\end{equation}

Release is local and causal.  Convertible momentum may return to translation
only while its owning contact remains geometrically meaningful and the current
constraint permits outward motion.  Released momentum changes velocity:

\begin{equation}
\Delta\mathbf v_i=
\frac{\Delta\mathbf p_{\mathrm{release}}}{m_i}.
\end{equation}

A nonzero remainder must never be silently erased.  On contact separation it
must be released, transferred to the opposing body or environment according to
an explicit rule, or reported as an unresolved model error.

\subsubsection{External Momentum}

Momentum transferred to a stationary wall belongs to the external
environment.  Momentum transferred to a piston belongs to the piston reaction
ledger.  These reaction ledgers are necessary when evaluating conservation of
the modeled system but do not require integrating wall position.

\subsection{Energy Scope}

The active model has no internal energy state and no energy-driven correction.
Kinetic or contact-potential energy may be calculated for diagnostics, but
those values do not determine force, containment, release, or position.

Momentum and energy are distinct quantities.  A force interaction transfers
impulse and may perform work, but missing momentum cannot be renamed heat.
Thermal energy, damping, plasticity, and irreversible dissipation require a
separate future material model.

\subsection{Multi-Contact Ownership}

One source may contact several particles and walls in the same frame.  Each
contact owns its identity, geometry, normal, target velocity, maximum-depth
state, and convertible momentum.  Source-level force is the vector sum of all
contact forces.  Source-level containment is one projection against the
complete active constraint set.

Sequentially correcting one pair and then another is not the approved model,
because the later correction can invalidate the earlier pair and introduces
order dependence.  Likewise, a single particle-level internal vector cannot
reliably represent several differently oriented contact ledgers.

\subsection{Lifecycle and Domain Rules}

The ordinary simple and reservoir models generate mobile particles active
from frame zero.  Runtime reservoir birth and deferred admission are retired.
Permanent life state is:

\begin{itemize}
    \item zero: active mobile particle;
    \item negative: dead mobile particle.
\end{itemize}

Two generator-scoped exceptions intentionally use positive \texttt{Data.w} as
a scheduled release frame:
\begin{itemize}
    \item \texttt{GenStreaming.py} uses positive \texttt{Data.w} for FreeStream
          inlet waves;
    \item \texttt{GenLightingData.py} uses positive \texttt{Data.w} for
          explicit lighting and visualization sequences.
\end{itemize}
In those records, the particle is pending until the current frame reaches the
stored value, then participates as an active mobile particle.  This is a
precomputed generator release schedule, not runtime birth admission, and it
does not write or relocate target particles.

After integration, a mobile particle outside the rectangular cell array is an
error.  A mobile particle outside any configured death plane becomes dead.
Death bounds never apply to boundary markers.

Dead state is not a particle type.  A dead mobile particle remains a mobile
record with negative \texttt{Data.w}.  Compute and vertex paths must both test
that lifecycle state before treating the record as a source, target,
cell-occupancy participant, or rendered particle.

The rectangular cell index is

\begin{equation}
i=x+W(y+Hz),
\end{equation}

with independent width \(W\), height \(H\), and depth \(L\).  Every axis is
validated against its own dimension.

\subsection{Python and GLSL Parity}

\texttt{ForceDynamics.py} is the shader-shaped model.  Its physical geometry,
contact force, wall evaluation, piston evaluation, maximum-depth constraints,
integration, and error rules define the GLSL behavior.  The Python base class
may load configuration, create structures, run direct scans, and collect
diagnostics, but those harness operations are not shader physics.

Python direct target scans and GLSL cell-occupancy scans may differ in broad
phase only.  Given the same local contacts and frame-start state, they must
produce the same geometry, force, constraint, accepted velocity, position, and
error result.  Generated GLSL must read current position and velocity buffers
and write only alternate source buffers.

\subsection{Diagnostics}

Diagnostics observe the model and must not alter it.  Useful quantities
include:

\begin{itemize}
    \item particle translational momentum;
    \item convertible internal momentum per contact;
    \item wall and piston reaction impulse;
    \item contact penetration and maximum-depth state;
    \item timestep-resolution margin;
    \item optional kinetic and contact-potential energy estimates.
\end{itemize}

Rendering color, particle numbers, presentation quality, and boundary drawing
are visualization state.  They must not affect contact detection, force,
constraints, integration, or lifecycle.

\subsection{Current Implementation Deviations}

The following behavior exists in the current Python implementation but is not
the approved final momentum model:

\begin{enumerate}
    \item ordinary compressive force impulse is copied into an internal
          momentum ledger even though it already changes translational
          velocity;
    \item maximum-depth blocked momentum is added to that same ledger;
    \item rebound drains the ledger diagnostically without applying the
          released momentum to particle velocity;
    \item residual storage is cleared when no contact remains;
    \item storage is aggregated on the particle rather than owned per contact.
\end{enumerate}

The required correction is to reserve convertible internal momentum solely for
blocked constraint momentum, store it per contact, apply physical release to
translation, and define explicit transfer or error behavior for every nonzero
remainder.  This correction follows physical-contact regression testing; it
must not reintroduce effective radii, suppressed overlap, pairwise target
writes, or order-dependent iterative correction.

\subsection{Open Model Work}

The remaining theory work is deliberately narrow:

\begin{itemize}
    \item define the exact per-contact convertible-momentum release rule;
    \item define reciprocal accounting for mobile targets without violating
          source ownership;
    \item define stationary-wall and piston transfer at contact termination;
    \item add swept-contact validation for complete between-frame crossings;
    \item generalize the two-dimensional constraint solver before enabling
          unrestricted three-dimensional multi-contact dynamics.
\end{itemize}
